\chapter{MCU、裸机与 FreeRTOS}
\label{chap:mcu-freertos}

微控制器（Microcontroller Unit, MCU）把处理器核、片上存储器、时钟、复位和常用外设集成在单一芯片中。
与运行完整 Linux 的应用处理器相比，MCU 系统通常具有更确定的启动路径、更有限的存储资源和更严格的实时、功耗约束。
本章以通用 32 位 MCU 为对象，说明从复位入口、裸机运行时到 FreeRTOS 多任务系统的关键机制。
FreeRTOS 的 API、配置项和移植层约束应以项目固定的内核版本及对应手册为准\cite{freertos-kernel-book,freertos-reference}；寄存器地址、中断号和时钟配置则必须以目标器件参考手册、数据手册和勘误表为准。

MCU 软件的正确性至少包含四个相互约束的维度：功能结果正确、截止期内完成、资源使用有界、故障后进入可解释状态。能够正常启动和运行示例负载，只证明了其中很小的一部分。

\section{MCU 系统模型}

典型 MCU 包含处理器核、片上 Flash、SRAM、中断控制器、DMA 控制器以及 GPIO、定时器、串行接口和模拟外设。
处理器通过存储器映射 I/O 访问外设寄存器；外设通过中断报告异步事件，也可以通过 DMA 在外设与存储器之间传输数据。

\subsection{存储器与地址空间}

MCU 固件通常由以下区域组成：
\begin{itemize}
  \item \textbf{代码和只读数据}：保存在非易失性存储器中，复位后可直接取指或复制到 RAM 执行；
  \item \textbf{已初始化数据段}：初值保存在镜像中，启动代码在进入 C 运行时前将其复制到 RAM；
  \item \textbf{未初始化数据段}：启动代码把 BSS 区清零；
  \item \textbf{堆与栈}：由链接脚本划定边界，必须结合最坏情况进行容量设计；
  \item \textbf{外设窗口}：用于访问控制、状态和数据寄存器，通常具有强顺序或设备内存属性。
\end{itemize}

链接脚本不是单纯的构建配置，它定义了固件镜像与芯片地址空间之间的契约。
工程至少应通过链接期断言或镜像检查确认代码、数据、启动向量、任务栈、持久化区和 Bootloader 保留区互不重叠，并记录镜像入口、加载地址、运行地址和初始栈边界。链接、重定位与启动运行时的详细关系见第~\ref{chap:compiler-linker-binary}~章。

\subsection{时钟、复位与启动模式}

复位释放后，处理器从架构定义或芯片配置指定的入口开始执行。
启动代码通常依次完成栈指针建立、异常向量安装、数据段初始化、BSS 清零、必要的时钟和存储器初始化，最后进入 \texttt{main()}。

时钟树决定处理器、总线和外设的实际工作频率。修改 PLL、分频器或时钟源时必须遵守电压、Flash 等待周期、锁定时间和切换顺序约束。
对于安全相关或高可靠系统，还应记录复位原因，区分上电复位、看门狗复位、欠压复位、软件复位和外部复位，并据此选择冷启动或受控恢复路径。

\subsection{复位后的状态分类}

复位不会自动把整个系统恢复为“全新状态”。复位域之外的外设、备份寄存器、retention SRAM、外部存储器和另一处理器核可能继续保留旧状态。启动代码应把状态分为：
\begin{itemize}
  \item \textbf{必须重新初始化}：时钟、IRQ 路由、DMA descriptor、驱动所有权和易失协议状态；
  \item \textbf{允许保留但必须校验}：崩溃记录、启动计数、校准参数和恢复事务；
  \item \textbf{必须在使用前清除}：可能包含密钥、旧指针或跨版本不兼容布局的内存；
  \item \textbf{由其他执行域拥有}：外部看门狗、协处理器、共享外设和通信对端，不能由本核单方面假定已复位。
\end{itemize}

初始化过程应具有阶段号或进度标记。若系统在初始化中途再次复位，下一次启动才能区分“稳态运行后故障”和“启动阶段反复失败”，并选择限次重试、回退镜像或安全模式。

\section{裸机软件结构}

裸机并不等同于“没有架构”。一个可维护的裸机系统通常仍应区分板级初始化、硬件抽象、设备驱动、协议与业务逻辑，并明确中断上下文和主循环上下文之间的接口。

\subsection{执行上下文契约}

每个可调用接口都应声明允许的执行上下文，而不是只声明参数类型：
\begin{itemize}
  \item \textbf{复位与早期启动}：C 运行时、全局对象、堆、时钟和中断可能尚不可用；
  \item \textbf{ISR}：不能阻塞，栈预算独立，且只能调用目标移植层明确允许的 API；
  \item \textbf{任务或主循环}：可以阻塞的接口必须同时定义超时、取消和资源回收语义；
  \item \textbf{内核回调}：软件定时器、Idle hook、内存失败 hook 和异常 hook 运行在不同上下文，不能相互套用约束；
  \item \textbf{Fault 与复位路径}：普通锁、日志服务和动态分配器可能已经损坏，必须使用最小依赖的证据保全路径。
\end{itemize}

接口文档至少应标注“可否阻塞、可否从 ISR 调用、是否改变 IRQ 屏蔽、缓冲区所有权、最长执行时间和失败后状态”。这类契约可以被静态检查、代码评审和测试共同验证。

\subsection{主循环与事件驱动}

最简单的固件采用轮询主循环。轮询便于理解，但当处理时间不可控或外设数量增加时，会造成响应延迟和无效功耗。
更常见的实现是由中断产生事件，主循环按优先级或状态机处理事件：

\begin{lstlisting}[language=C]
#include <stdatomic.h>

static _Atomic uint32_t pending_events;

for (;;) {
    uint32_t events = atomic_exchange_explicit(
        &pending_events, 0U, memory_order_acq_rel);

    if ((events & EVENT_RX_READY) != 0U)
        protocol_process_rx();
    if ((events & EVENT_SAMPLE_READY) != 0U)
        control_process_sample();
    if ((events & EVENT_FAULT) != 0U)
        system_enter_safe_state();

    if (events == 0U)
        platform_idle_if_no_events(&pending_events);
}
\end{lstlisting}

事件位只能表达“某类事件已经发生”，不能保留重复次数或携带大量数据。
需要无丢失传递时，应使用计数器、环形缓冲区或固定容量消息队列，并定义队列满时的策略。
ISR 可以用 \texttt{atomic\_fetch\_or\_explicit(..., memory\_order\_release)} 发布事件，但前提是目标上的该原子操作确实 lock-free。若编译器生成库调用或内部锁，它就不适合异常上下文；工程应检查生成代码或 \texttt{atomic\_is\_lock\_free()}，否则使用经过验证的短临界区或架构专用原语。
\texttt{platform\_idle\_if\_no\_events()} 必须在受控 IRQ 屏蔽范围内重新检查事件，再以平台保证不会丢失唤醒的顺序执行睡眠；直接在“检查为空”后调用 WFI 会留下检查与休眠之间的竞态窗口。

\subsection{中断设计}

中断服务程序（Interrupt Service Routine, ISR）的首要目标是有界、可分析地完成紧急处理。
通常只读取或确认硬件状态、保存时间敏感数据并通知后续处理，不在 ISR 中执行阻塞操作、复杂协议解析或不可预测的动态内存分配。

设计中断时应明确：
\begin{itemize}
  \item 中断源的触发方式、清除顺序和重复触发条件；
  \item 优先级分组、抢占关系和最大嵌套深度；
  \item ISR 与线程或主循环共享对象的原子性与内存可见性；
  \item 最坏中断响应时间和最长关中断时间；
  \item 中断风暴、外设失效和队列溢出时的降级措施。
\end{itemize}

中断延迟必须按完整链路测量：硬件事件到 IRQ 置位、优先级仲裁、被更高优先级 ISR 或临界区阻塞、入口保存、驱动确认、任务唤醒以及调度切换。只测 ISR 函数体会漏掉最常见的尾部延迟。对于电平触发中断，还要验证处理完成前外设状态已真正解除，否则退出 ISR 后会立即再次进入。

\subsection{DMA 与缓存一致性}

DMA 可以降低 CPU 搬运负担，但会引入缓冲区所有权和一致性问题。
在没有数据缓存的 MCU 上，仍需防止 CPU 与 DMA 同时修改同一缓冲区；在带缓存的 MCU 或 SoC 上，还必须按照平台规则执行缓存清理、失效或使用一致性内存。
DMA 缓冲区还应满足地址对齐、可访问存储区、传输长度和外设突发边界等要求。

推荐把每个 descriptor 或缓冲区建模为显式状态机，例如 \texttt{CPU\_FREE -> CPU\_READY -> DMA\_OWNED -> CPU\_DONE}。所有权交接点同时承担长度、代次号和内存可见性检查；超时取消时，必须先证明 DMA 已停止访问，才能回收或复用缓冲区。双缓冲只能吸收有限抖动，生产者长期快于消费者时仍需要丢弃、背压或降采样策略。

\section{实时系统基础}

实时系统的正确性同时取决于计算结果和结果产生的时间。
“运行得快”不等于实时；实时设计要求关键路径的最坏执行时间、阻塞时间和调度延迟能够被约束，并在规定截止期内完成。
本章侧重把需求映射到 MCU 与 FreeRTOS 机制；响应时间分析、release jitter、阻塞项和语言内存模型见第~\ref{chap:realtime-language-memory-model}~章。

\subsection{实时性分类}

\begin{description}
  \item[硬实时] 任何关键截止期违约都可能造成不可接受的系统后果；
  \item[固实时] 偶发超期结果通常已经失去价值，但不一定立即导致灾难性后果；
  \item[软实时] 延迟主要影响服务质量，可以用统计分位数描述。
\end{description}

实时需求应写成可测量指标，例如“外部中断到执行器更新的最坏延迟不超过 $200\,\mu s$”，而不是笼统地写成“响应迅速”。

\subsection{任务与端到端链路}

单个任务至少需要定义周期或最小到达间隔、截止期、最坏执行时间、最大阻塞时间和允许抖动。端到端链路还要覆盖 IRQ、DMA、队列、任务切换和输出外设，而不能把每个局部函数都满足预算误认为整条链路满足预算。

同一功能在冷启动、正常运行、校准、升级和故障恢复模式下可能具有不同执行时间与任务集合。模式切换验收要覆盖旧模式工作尚未排空而新模式任务已经释放的重叠窗口。过载策略也必须预先定义：保留最新采样、丢弃低优先级工作、降低算法质量、进入安全状态，还是复位恢复。

\subsection{调度与优先级}

FreeRTOS 常用配置采用基于优先级的抢占式调度，也可以按配置使用协作式调度和同优先级时间片\cite{freertos-kernel-book}。
高优先级任务只应承担更紧迫的工作，不能用优先级表达模块重要程度。
如果一个高优先级任务长期处于就绪态，所有较低优先级任务都可能饥饿。

优先级分配应有任务模型依据。周期越短不总是越紧急，设备服务任务、共享资源服务端和故障处理路径还会产生优先级依赖。多个同优先级任务启用时间片后，单次响应会受到同级任务数量和 tick 相位影响；硬实时任务更适合使用可分析的不同优先级或明确的协作调度点。

周期任务应优先使用绝对周期延时，避免把任务执行时间累积到周期误差中：

\begin{lstlisting}[language=C]
static void control_task(void *argument)
{
    TickType_t release = xTaskGetTickCount();

    for (;;) {
        acquire_inputs();
        run_control_loop();
        update_outputs();
        xTaskDelayUntil(&release, pdMS_TO_TICKS(10));
    }
}
\end{lstlisting}

调度节拍的分辨率不是系统真实时间精度。亚节拍级时间戳、脉冲捕获或精确输出应由硬件定时器完成，而不是依赖任务延时。

大多数 MCU FreeRTOS port 是单核调度器；若选择支持 SMP 的内核与移植层，还必须重新审查任务亲和性、跨核唤醒、全局临界区、共享 cache 和每核中断路由，不能直接复用单核最坏情况结论。

\section{FreeRTOS 内核对象}

FreeRTOS 内核提供任务、队列、信号量、互斥量、事件组、流缓冲区、消息缓冲区、软件定时器和任务通知等基本机制。
对象选择应以通信语义为依据，而不是以 API 熟悉程度为依据；边界行为还要核对所选内核版本的参考手册\cite{freertos-reference}。

常见映射包括：
\begin{description}
  \item[最新状态] 长度为 1 的队列或显式 mailbox，允许旧值被覆盖，但必须记录覆盖次数；
  \item[每次事件] 多元素队列或计数语义，必须为满队列定义丢弃、背压或故障路径；
  \item[一对一唤醒] 任务通知，适合计数、位标志或简单值，不自动携带复杂对象生命周期；
  \item[字节或消息流] Stream Buffer 或 Message Buffer，默认单写者、单读者假设及多方访问限制应按版本验证；
  \item[组合条件] Event Group，位标志会合并事件，不适合表达每次发生次数；
  \item[共享资源所有权] 互斥量或专用服务任务，而不是把二值信号量当作所有锁的替代品。
\end{description}

任何有界对象都需要容量依据。平均输入速率低于平均处理速率并不足够，仍要覆盖允许的最大 burst、消费者最长阻塞窗口和恢复时间。

\subsection{任务与栈}

每个任务具有独立栈、任务控制块、优先级和运行状态。任务栈至少需要容纳最深调用链、局部对象、保存的上下文以及库函数的额外需求。
栈容量应通过静态分析、压力场景下的高水位测量和溢出检测共同验证，不能只根据正常运行结果估算。

FreeRTOS API 中的栈深度和高水位通常以 \texttt{StackType\_t} 元素而不是字节表示，具体类型取决于移植层。高水位只反映已经执行过的路径；格式化、浮点、异常嵌套、库回调和罕见错误分支可能显著增加栈需求。还要分别预算任务栈与异常使用的主栈，防止只监控任务栈而遗漏 ISR 嵌套溢出。

任务函数通常不应返回；需要删除任务时，应先释放由任务独占的资源，再调用 \texttt{vTaskDelete()}。
删除另一个仍在运行或可能持锁的任务会破坏对象不变量，应优先通过停止请求、确认退出和统一回收实现受控生命周期。对于长期运行的产品，优先采用静态创建或启动阶段一次性创建，减少堆碎片和运行时资源耗尽的不确定性。

\subsection{队列与任务通知}

队列按值复制固定大小的元素，适合传递小型命令、采样结果或指针所有权。
向队列发送指针时必须定义所指对象的生命周期、所有者和回收责任。

任务通知直接存放在任务控制块中，开销通常低于单独创建信号量或队列，适合一对一事件、计数或位标志通知。
如果多个生产者需要传递不同类型的数据，或者需要独立排队和背压，队列通常更清晰。

通知的 overwrite、increment、set-bits 和 no-overwrite 模式具有不同丢失语义。接收方在清除位或取出计数前后发生的新通知是否会被合并，必须通过 API 语义和并发测试确认，不能把所有通知都理解为“可靠消息”。

\subsection{互斥量、信号量与优先级反转}

互斥量用于保护具有所有权概念的共享资源，并支持优先级继承；二值或计数信号量用于表示事件或可用资源数量。
ISR 不能获取可能阻塞的互斥量。

当低优先级任务持有高优先级任务所需的锁，而中优先级任务持续抢占低优先级任务时，会发生无界优先级反转。
FreeRTOS 互斥量的优先级继承能够缓解典型场景，但不能替代良好的锁设计。临界区应保持短小，避免在持锁期间进行阻塞 I/O、长时间计算或调用未知时延的回调。
优先级继承不能消除循环等待，也不能自动证明嵌套多把锁时的阻塞上界。需要共享串行外设时，采用单一服务任务并通过队列提交事务，通常比让多个任务持锁执行完整设备事务更容易分析。

\subsection{中断安全 API}

FreeRTOS 为 ISR 提供带 \texttt{FromISR} 后缀的 API。ISR 唤醒更高优先级任务后，应按照移植层规定请求一次上下文切换：

\begin{lstlisting}[language=C]
void DMA_IRQHandler(void)
{
    BaseType_t higher_priority_woken = pdFALSE;

    dma_acknowledge_interrupt();
    vTaskNotifyGiveFromISR(data_task_handle, &higher_priority_woken);
    portYIELD_FROM_ISR(higher_priority_woken);
}
\end{lstlisting}

能够调用 FreeRTOS API 的中断优先级范围由端口和 \texttt{FreeRTOSConfig.h} 共同约束。
在采用数值越小优先级越高的中断控制器上，配置错误尤其容易造成误解。
工程必须按照目标端口文档核对 \texttt{configMAX\_SYSCALL\_INTERRUPT\_PRIORITY} 等参数，并启用内核断言发现非法调用。

在典型 Cortex-M port 中，高于系统调用阈值的 IRQ 可以抢占内核临界区，因此绝不能调用内核 API；它只能操作独立的硬件状态或经过证明的无锁通道，再由较低优先级 IRQ 或任务完成递交。配置值还可能分为“未移位的逻辑优先级”和“写入 NVIC 字段的编码值”，应由一个集中配置模块生成并用编译期断言核对 \texttt{\_\_NVIC\_PRIO\_BITS}，避免各驱动自行移位。

\section{时间、内存与低功耗}

\subsection{Tick、超时与回绕}

内核 tick 是调度时间基准，不一定是墙上时间，也不保证与外设时间戳同源。将毫秒转换为 tick 时要检查舍入为零和整数溢出；需要至少阻塞一个 tick 的路径应显式处理最小值，而长超时不应假定 \texttt{pdMS\_TO\_TICKS()} 在所有类型宽度上都不会溢出。

tick 计数器最终会回绕。周期延时和复杂超时优先使用内核提供的相对时间 API、\texttt{TimeOut\_t} 与 \texttt{xTaskCheckForTimeOut()} 等版本匹配机制，不要用 \texttt{now >= deadline} 直接比较跨回绕的绝对值。无限等待也要谨慎：若资源永久不再到达，任务、锁和看门狗健康判定都需要可观测的超时或取消路径。

\subsection{软件定时器}

软件定时器回调在定时器服务任务中串行执行，不属于硬件中断。
回调不得长时间阻塞，否则会推迟其他定时器命令和回调。需要较长处理时，回调应只通知工作任务。
创建、启动、停止和复位定时器通常通过定时器命令队列递交；队列满、daemon 任务优先级不足或回调自发产生过多命令都会形成延迟。验收应同时观察回调最迟时间和命令队列峰值，而不是只检查平均周期。

\subsection{动态内存策略}

FreeRTOS 提供多种示例堆实现，但它们具有不同的分配、释放、合并和确定性特征。
安全或硬实时系统应明确是否允许调度器启动后继续动态分配，并为分配失败设置可验证的处理路径。
无论采用静态还是动态分配，都应统计峰值使用量并保留设计裕量。

静态分配也不等于自动安全：控制块、栈和队列存储区的生命周期必须覆盖内核对象，地址和对齐必须满足端口要求，禁止把局部数组作为长期对象后离开作用域。只启用静态分配时，内核仍可能要求应用提供 Idle task 以及启用软件定时器时的 Timer task 内存回调；函数签名和所需对象必须与固定内核版本一致。

资源预算应形成可核对的清单：每个任务的栈、所有静态队列与缓冲区、内核对象控制块、ISR 主栈、DMA 专用区和启动余量。堆最小剩余量、栈高水位和队列峰值是运行证据，但不能替代对尚未覆盖路径的静态预算。

\subsection{Tickless Idle}

当一段时间内没有任务需要运行时，Tickless Idle 可以停止周期节拍并进入低功耗状态。
其正确实现依赖低功耗定时源、唤醒源、时钟恢复时间和移植层补偿。
系统必须验证睡眠前后的时间连续性、外设状态、调试接口影响以及最短可获益睡眠时间。

“检查无任务可运行”到真正执行睡眠指令之间存在竞态。移植层必须在受控 IRQ 屏蔽范围内再次确认睡眠条件，并保证已经到达的唤醒事件不会丢失。唤醒后还要按顺序恢复时钟、电压、Flash 等待周期和外设，再允许依赖这些资源的 ISR 或任务运行。更完整的系统级低功耗契约见第~\ref{chap:low-power-system}~章。

\section{FreeRTOS 工程配置}

\texttt{FreeRTOSConfig.h} 是应用与内核之间的重要配置接口。至少应审查：
\begin{itemize}
  \item 固定的内核 tag 或 commit、port 路径、补丁集和许可证信息；
  \item CPU 时钟、tick 频率和最大优先级；
  \item 抢占、时间片、低功耗和软件定时器配置；
  \item 静态与动态内存支持；
  \item 栈溢出、分配失败、断言和运行时统计钩子；
  \item 中断优先级位宽及可调用内核 API 的边界；
  \item 启用的 API、兼容宏和对象数量，避免无意扩大资源占用或改变语义；
  \item tick 类型宽度、运行时计数器宽度及其回绕周期。
\end{itemize}

内核配置属于系统架构的一部分，应纳入代码评审和版本管理，不能由单个示例工程直接复制后长期不审查。
建议在产物 manifest 中记录内核版本、\texttt{FreeRTOSConfig.h} 摘要、编译器与优化等级，并通过编译期断言检查时钟、优先级范围、类型宽度和必须启用的 hook。升级内核或更换 port 时，应重新执行调度、ISR、栈、低功耗和故障注入测试，而不是只依赖 API 编译兼容。

\section{最小可运行 FreeRTOS 工程}

一个适合验证移植层的最小工程应同时覆盖任务切换、阻塞队列、ISR 唤醒和周期调度，而不是只创建一个永不阻塞的闪灯任务。建议目录如下：

\begin{lstlisting}[language=bash]
freertos-minimal/
  CMakeLists.txt
  linker.ld
  startup.S
  include/FreeRTOSConfig.h
  src/main.c
  src/board.c
  src/freertos_hooks.c
  src/interrupts.c
  third_party/FreeRTOS-Kernel/
\end{lstlisting}

示例由按键 ISR 产生事件，控制任务切换 LED 模式，周期任务按照绝对周期刷新输出。队列、任务、栈和内核 Idle task 内存均静态分配，从而使调度器启动后的内存需求固定。

\begin{lstlisting}[language=C,caption={静态对象和任务入口}]
#include "FreeRTOS.h"
#include "queue.h"
#include "task.h"
#include "board.h"
#include <stdatomic.h>

typedef enum { LED_SLOW, LED_FAST } led_mode_t;
static StaticQueue_t mode_queue_cb;
static uint8_t mode_queue_storage[sizeof(led_mode_t)];
static QueueHandle_t mode_queue;

static StaticTask_t control_tcb, led_tcb;
static StackType_t control_stack[256], led_stack[256];
static _Atomic led_mode_t current_mode = LED_SLOW;

static void control_task(void *arg)
{
    (void)arg;
    led_mode_t mode;

    /* Clear peripheral/NVIC pending state before unmasking the IRQ. */
    board_button_irq_start();

    for (;;) {
        xQueueReceive(mode_queue, &mode, portMAX_DELAY);
        atomic_store_explicit(&current_mode, mode, memory_order_release);
    }
}

static void led_task(void *arg)
{
    (void)arg;
    TickType_t release = xTaskGetTickCount();

    for (;;) {
        board_led_toggle();
        led_mode_t mode = atomic_load_explicit(&current_mode,
                                               memory_order_acquire);
        TickType_t period = mode == LED_FAST ?
            pdMS_TO_TICKS(100) : pdMS_TO_TICKS(500);
        xTaskDelayUntil(&release, period);
    }
}
\end{lstlisting}

\begin{lstlisting}[language=C,caption={系统初始化与 ISR 事件投递}]
int main(void)
{
    board_init();
    configASSERT(atomic_is_lock_free(&current_mode));

    mode_queue = xQueueCreateStatic(1, sizeof(led_mode_t),
        mode_queue_storage, &mode_queue_cb);
    configASSERT(mode_queue != NULL);

    configASSERT(xTaskCreateStatic(control_task, "control", 256, NULL,
        2, control_stack, &control_tcb) != NULL);
    configASSERT(xTaskCreateStatic(led_task, "led", 256, NULL,
        1, led_stack, &led_tcb) != NULL);

    vTaskStartScheduler();
    board_fault_stop();
}

void BUTTON_IRQHandler(void)
{
    static led_mode_t next = LED_FAST;
    BaseType_t wake = pdFALSE;

    board_button_irq_ack();
    xQueueOverwriteFromISR(mode_queue, &next, &wake);
    next = next == LED_FAST ? LED_SLOW : LED_FAST;
    portYIELD_FROM_ISR(wake);
}
\end{lstlisting}

\texttt{xQueueOverwriteFromISR()} 只适用于长度为 1 的队列，因此这里的对象语义是“最新模式邮箱”，而不是保留每一次按键的事件队列。若业务要求不丢失每次边沿，应改用多元素队列和 \texttt{xQueueSendFromISR()}，并定义满队列处理策略。
IRQ 由控制任务在调度器开始运行后启用，避免启动前的 ISR 请求一次尚不可执行的任务切换；\texttt{board\_button\_irq\_start()} 还应按芯片规则清除外设标志与 NVIC pending。

关闭动态分配时，应用必须为 Idle task 提供静态内存。下列签名适用于使用 \texttt{configSTACK\_DEPTH\_TYPE} 的内核版本；固定其他版本时应按其头文件调整第三个参数类型：

\begin{lstlisting}[language=C,caption={静态 Idle task 内存回调}]
static StaticTask_t idle_tcb;
static StackType_t idle_stack[configMINIMAL_STACK_SIZE];

void vApplicationGetIdleTaskMemory(
    StaticTask_t **tcb,
    StackType_t **stack,
    configSTACK_DEPTH_TYPE *depth)
{
    *tcb = &idle_tcb;
    *stack = idle_stack;
    *depth = configMINIMAL_STACK_SIZE;
}
\end{lstlisting}

若启用软件定时器，还要按同一版本接口提供 Timer service task 的静态内存。hook 中使用的所有对象必须具有静态生命周期。

关键配置至少应启用抢占、静态分配、断言和栈溢出检查，并关闭未使用的动态分配：

\begin{lstlisting}[language=C,caption={最小 FreeRTOSConfig.h 关键项}]
#define configUSE_PREEMPTION                 1
#define configUSE_TIME_SLICING                0
#define configUSE_TICKLESS_IDLE               0
#define configSUPPORT_STATIC_ALLOCATION      1
#define configSUPPORT_DYNAMIC_ALLOCATION     0
#define configUSE_TIMERS                      0
#define configCHECK_FOR_STACK_OVERFLOW       2
#define configUSE_MALLOC_FAILED_HOOK         0
#define configTICK_RATE_HZ                    1000
#define configMAX_PRIORITIES                  5
#define configMINIMAL_STACK_SIZE              PLATFORM_IDLE_STACK_WORDS
#define INCLUDE_vTaskDelay                    1
#define INCLUDE_xTaskDelayUntil               1
#define configASSERT(x) do { if (!(x)) board_fault_stop(); } while (0)
\end{lstlisting}

该片段只表达本例策略，不是完整配置文件；CPU 时钟、中断优先级、port 必需宏和类型配置必须来自目标移植层。工程的 CMake 目标应显式包含 FreeRTOS 内核、目标架构 port、heap 实现（仅在启用动态分配时）、启动文件和链接脚本。验收时必须观察 LED 周期、连续按键、队列覆盖语义、任务栈高水位、tick 回绕和非法中断优先级断言；仅能编译链接不代表移植正确。

\section{健康监控、看门狗与故障恢复}

硬件看门狗只能证明“某段代码仍能周期执行”，不能自动证明采样、控制、通信和输出链路都在前进。更可靠的做法是由健康监控任务收集各关键任务的进度代次、最后成功时间和错误状态，只有所有必需条件在窗口内满足时才喂狗。

健康条件应基于业务进展而不是任务是否被调度。例如通信任务循环运行但一直重试同一事务、控制任务周期唤醒但输入时间戳不再更新，都属于失去进展。监控任务自身不应具有绕过检查直接喂狗的后门；窗口看门狗还可检测过早喂狗和失控快速循环。

故障路径至少要区分：
\begin{itemize}
  \item \textbf{局部可恢复}：重置单一外设、重建队列或重新同步协议；
  \item \textbf{功能降级}：关闭非关键负载，保持安全控制与诊断；
  \item \textbf{受控系统复位}：先冻结执行器或进入安全状态，再保全最小证据并触发复位；
  \item \textbf{启动回退}：连续启动失败达到阈值后选择已知良好镜像或维护模式。
\end{itemize}

崩溃记录应包含复位原因、固件标识、任务或异常上下文、栈边界、关键队列水位、最近进度代次和单调时间。记录写入必须有界，并考虑掉电中断和 retention 区损坏；更完整的证据链与符号化方法见第~\ref{chap:debug-reliability}~章。

\section{架构与验证方法}

推荐把系统划分为中断层、驱动层、服务层和应用层。中断层只完成硬件事件收敛；驱动层管理外设状态和缓冲区；服务层提供协议、存储或控制能力；应用层实现产品状态机。
跨层接口应显式描述执行上下文、阻塞属性、超时、缓冲区所有权和错误返回。

FreeRTOS 系统的验证至少包括：
\begin{enumerate}
  \item 固定内核、port、配置、工具链、镜像和硬件版本，形成可复现实验 manifest；
  \item 测量关键 ISR、任务、临界区和端到端链路的最长执行与响应时间；
  \item 在最大 burst 和消费者最长阻塞下验证队列深度、背压和溢出处理；
  \item 覆盖罕见错误路径后检查所有任务栈、ISR 主栈、堆和静态池峰值；
  \item 验证 tick 回绕、长超时、连续超时重试、软件定时器拥塞和时间源漂移；
  \item 注入中断风暴、DMA 超时、总线错误、传感器失效和内存分配失败；
  \item 验证看门狗只能由健康监控路径喂狗，并能识别“任务运行但业务无进展”；
  \item 在任意初始化阶段复位，验证故障记录、持久化数据、安全状态和启动回退；
  \item 在发布构建、目标优化等级、最低电压、最高温度和低功耗切换下重复测试。
\end{enumerate}

\newpage
\section{设计检查清单}

\begin{itemize}
  \item 内核版本、port、补丁、配置摘要、工具链与硬件版本是否进入产物 manifest？
  \item 是否为每个任务定义了周期或触发条件、优先级依据、截止期和最坏执行时间？
  \item 每个接口是否标注可调用上下文、阻塞性、超时、所有权和最长执行时间？
  \item ISR 是否有明确的最大执行时间，且优先级允许其调用所用的中断安全 API？
  \item 队列满、通知合并、软件定时器命令拥塞、缓冲区耗尽和锁超时是否都有确定策略？
  \item 是否验证了任务栈、ISR 主栈、堆、静态对象、DMA 缓冲区和持久化区域的边界？
  \item 中断优先级与 FreeRTOS 系统调用优先级配置是否经过目标板验证？
  \item tick 回绕、毫秒转换、无限等待、取消和低功耗时间补偿是否有测试证据？
  \item 看门狗是否检查业务进展，而不是仅由某个周期任务无条件喂狗？
  \item 系统在外设失效、任务卡死、初始化中途复位、欠压和看门狗复位后是否进入可解释状态？
  \item 测试是否覆盖最坏 burst、最低电压、最高温度、发布优化和低功耗切换？
\end{itemize}

第~\ref{chap:case-dvrapp-mcu}~章进一步说明如何把多类 Cortex-M 芯片、多套工具链、IAP、DMA UART、传感器和 MCU--CPU 协议组织为可迁移的 FreeRTOS 工程。

第~\ref{chap:sampling-control-dsp}~章进一步讨论采样、滤波、定点数值与控制任务的时序约束。
